\section{Описание методов}
\subsection{Статистическое и транспортное согласование}
Пусть $z$ --- вектор признаков, извлечённый из пространственной позиции карты сети, а $P_s$ и $P_t$ --- распределения исходных и целевых признаков. Ядерный критерий MMD сравнивает их средние вложения в пространство $\mathcal H$~\cite{mmd}:
\begin{equation}
 L_{\mathrm{MMD}}=\left\|\mathbb E_{P_s}\phi(z)-\mathbb E_{P_t}\phi(z)\right\|_{\mathcal H}^2.
\end{equation}
На практике используются значения ядра $k(z,z')=\langle\phi(z),\phi(z')\rangle$. Сеть обновляется по исходной ошибке и оценке MMD, вычисленной на выбранных признаках двух доменов. Выбор ядра, его масштаба и способа выборки пикселей влияет на результат. Полная попарная оценка требует квадратичного числа операций по числу выбранных признаков.

Deep CORAL согласует ковариации $C_s$ и $C_t$ признаков размерности $d$~\cite{coral}:
\begin{equation}
 L_{\mathrm{CORAL}}=\frac{1}{4d^2}\|C_s-C_t\|_F^2.
\end{equation}
Ковариации вычисляются после центрирования признаков. Их совпадение не гарантирует совпадения средних или полных распределений. В сегментационной модели оба критерия можно применять к выборке пространственных признаков; согласование без учёта классов способно сблизить семантически разные объекты.

Оптимальный транспорт задаёт матрицу неотрицательных переносимых масс $\gamma$. Для исходных весов $a$ и целевых весов $b$, каждый из которых суммируется в единицу, рассматривается задача
\begin{equation}
 \min_{\gamma\in\Pi(a,b)}\;
 \langle\gamma,M\rangle+\varepsilon\sum_{i,j}\gamma_{ij}(\log\gamma_{ij}-1),
\end{equation}
где $\Pi(a,b)=\{\gamma\geq0:\gamma\mathbf1=a,\gamma^T\mathbf1=b\}$, $M$ --- стоимость сопоставления, $\varepsilon>0$. Принимается $0\log0=0$. Алгоритм Синкхорна попеременно нормирует строки и столбцы ядра $K_{ij}=\exp(-M_{ij}/\varepsilon)$~\cite{sinkhorn}. После расчёта плана обновляются параметры сети. Семантическая информация может включаться в стоимость; память для матрицы попарных стоимостей ограничивает число используемых признаков.

\subsection{Состязательное согласование}
Дискриминатор $D_\varphi$ оценивает вероятность принадлежности исходному домену по признакам либо предсказаниям $r_\theta(x)$. Базовая постановка имеет вид
\begin{equation}
 \begin{split}
 V(\theta,\varphi)={}&\mathbb E_s\log D_\varphi(r_\theta(x))\\
 &+\mathbb E_t\log[1-D_\varphi(r_\theta(x))],\\
 &\min_\theta\max_\varphi\;L_s(\theta)+\lambda V(\theta,\varphi).
 \end{split}
\end{equation}
Максимизация обучает различать домены, минимизация по $\theta$ затрудняет это различение. Используются попеременные градиентные шаги; конкретные алгоритмы могут менять генераторную потерю для устойчивости. AdaptSegNet применяет этот принцип к картам классов~\cite{tsai}. Основные ограничения --- неустойчивость минимаксного обучения и возможность подавления полезных различий. Низкая точность дискриминатора сама по себе не подтверждает правильность сегментации.

\subsection{Минимизация энтропии}
Для вероятностей классов $p_c$ нормированная энтропия равна
\begin{equation}
 H(p)=-\frac1{\log C}\sum_{c=1}^{C}p_c\log p_c.
\end{equation}
Её среднее по целевым пикселям добавляется к $L_s$. Критерий стимулирует уверенные ответы, но уверенная ошибка также имеет малую энтропию. Поэтому важно контролировать вес регуляризации и распределение предсказываемых классов. В ADVENT различаются прямая минимизация энтропии MinEnt и состязательная обработка карт самоинформации AdvEnt~\cite{advent}; их нельзя считать одной и той же процедурой оптимизации.

\subsection{Самообучение с псевдоразметкой}
Модель с фиксированными на текущем шаге параметрами $\bar\theta$ формирует метки $\hat y_u=\arg\max_c p_{\bar\theta,c}(x^t)_u$. Для отбора можно использовать маску $w_u=\mathbf1[\max_c p_{\bar\theta,c}(x^t)_u\geq\tau_{\hat y_u}]$ и потерю
\begin{equation}
 L_{\mathrm{pl}}=-\frac{\sum_u w_u\log p_{\theta,\hat y_u}(x^t)_u}
 {\max(1,\sum_u w_u)}.
\end{equation}
Псевдометки и маска фиксируются при обновлении ученика. Затем предсказания уточняются, и обучение повторяется. Пороги могут зависеть от класса, поскольку единый высокий порог часто отбирает преимущественно простые и частые категории. Это обобщённая схема отбора, а не точное воспроизведение DACS или DAFormer.

В DACS части исходных изображений вместе с истинными метками смешиваются с целевыми областями и псевдометками~\cite{dacs}. DAFormer использует выборку исходных изображений с редкими классами~\cite{daformer}. Эти механизмы помогают сохранять классовое разнообразие, но не устраняют риск закрепления ошибок автоматически. Сравнение в табл.~\ref{tab:dacs} показывает, почему наличие псевдометок само по себе не гарантирует улучшения.

\subsection{Согласованность предсказаний}
Для двух допустимых преобразований одного изображения минимизируется расхождение ответов, например дивергенция Кульбака--Лейблера между ответом учителя $q_u$ и ученика $p_u$. Геометрические преобразования требуют предварительного совмещения координат. Параметры учителя могут обновляться экспоненциальным усреднением:
\begin{equation}
 \bar\theta_k=\alpha\bar\theta_{k-1}+(1-\alpha)\theta_k,\qquad 0\leq\alpha<1.
\end{equation}
Градиент через целевой ответ учителя не проводится. В MIC применяется перекрёстная энтропия между предсказаниями ученика по маскированному изображению и псевдометками учителя по полному изображению~\cite{mic}. Учитель сглаживает изменения параметров, но не получает независимой правильной разметки. Два одинаковых ответа могут быть одновременно ошибочными.

\subsection{Контрастивные критерии}
Контрастивное обучение сближает признаки одного класса и разделяет признаки разных классов. Положительные соответствия задаются исходными метками или целевыми псевдометками; отрицательные --- другими классами. Вариант с прототипами сравнивает признак с представительными векторами классов через softmax от сходства, делённого на температуру. Критерии типа InfoNCE задают относительное предпочтение положительных соответствий отрицательным.

В SePiCo классы представляются прототипами, банками признаков или распределениями~\cite{sepico}. Для исследования важны надёжность псевдометок, обновление прототипов, температура и состав отрицательных примеров. Ошибочная метка создаёт неверное семантическое соответствие. Поэтому контрастивный член следует оценивать как дополнение к фиксированной базовой схеме, отдельно проверяя его влияние на редкие классы.

\subsection{Проект сравнительного исследования}
Будущее сравнение предполагает единые разбиения данных, набор из 19 классов, начальные веса, разрешение входа и бюджет обучения. Сначала оценивается source only. Затем в той же модели по отдельности включаются выбранные критерии. Комбинации исследуются через исключение или добавление одного компонента при сохранении остальных условий. Полные реализации AdaptSegNet, DACS и DAFormer рассматриваются отдельно как ориентиры из литературы.

Для каждого запуска следует сохранять конфигурацию, случайное зерно, версию реализации, правило выбора контрольной точки, mIoU и IoU классов, время обучения и пиковую память. Если ресурсы позволяют, выполняются повторные запуски и сообщаются среднее и стандартное отклонение. Разности малой величины без оценки вариативности не считаются достаточным основанием для вывода о превосходстве метода. Снижение метрики относительно той же исходной модели фиксируется как отрицательный перенос.

Предстоит проверить, какие критерии дают прирост при одинаковых условиях и сохраняется ли он для редких классов. Набор вариантов и бюджет запусков зависят от вычислительных ресурсов. Обучение по этому протоколу ещё не проводилось.

